% AUTO-GENERATED. Edit the Markdown sources or notes.config.json.
% Sources: 02-Economics/2026-Macroeconomics/2026-09-07-Note.md

\section{The Science of Macroeconomics / 宏观经济学的科学方法}

\subsection{The Subject Matter of Macroeconomics / 宏观经济学的研究对象}

\textbf{Definition}

\begin{quote}
\textbf{Macroeconomics} is the study of the economy as a whole.
\end{quote}

宏观经济学研究整体经济的运行。典型问题包括：

\begin{itemize}
\tightlist
\item
  经济衰退为何发生，政府刺激政策为何可能有效；
\item
  房地产市场的问题如何传导至其他部门；
\item
  政府预算赤字如何影响劳动者、消费者、企业与纳税人；
\item
  通货膨胀为何提高生活成本；
\item
  国家间长期收入差异如何形成，哪些政策有助于经济增长；
\item
  贸易赤字如何影响国民福利。
\end{itemize}

三组核心宏观数据为：

\begin{longtable}[]{@{}ll@{}}
\toprule\noalign{}
指标 & 主要信息 \\
\midrule\noalign{}
\endhead
\bottomrule\noalign{}
\endlastfoot
实际人均 GDP（real GDP per capita） & 长期增长与短期波动 \\
通货膨胀率（inflation rate） & 生活成本的变化 \\
失业率（unemployment rate） & 劳动力市场状况与经济周期 \\
\end{longtable}

\subsection{Economic Models / 经济模型}

\textbf{Definition}

\begin{quote}
\textbf{Economic models} are simplified versions of a more complex reality. Irrelevant details are stripped away.
\end{quote}

模型用于描述变量关系、解释经济行为并分析政策后果。

\subsubsection{Endogenous and Exogenous Variables / 内生变量与外生变量}

\textbf{Definition}

\begin{quote}
The values of \textbf{endogenous variables} are determined in the model.

The values of \textbf{exogenous variables} are determined outside the model. The model takes their values and behavior as given.
\end{quote}

通常可以把模型理解为
\[
\text{外生变量与参数}
\longrightarrow
\text{模型机制}
\longrightarrow
\text{内生变量}.
\]

\subsubsection{Why We Need Multiple Models / 为什么需要多个模型}

不存在能够回答所有问题的单一模型。学习每个模型时，应明确：

\begin{enumerate}
\def\labelenumi{\arabic{enumi}.}
\tightlist
\item
  模型的假设；
\item
  哪些变量是内生变量，哪些是外生变量；
\item
  模型能够回答哪些问题，又忽略了哪些机制。
\end{enumerate}

\subsection{Flexible and Sticky Prices / 价格伸缩性与价格黏性}

\textbf{Definition}

\begin{quote}
\textbf{Market clearing} is an assumption that prices are flexible and adjust to equate supply and demand.
\end{quote}

该假设在不同时间尺度上的适用性不同：

\begin{longtable}[]{@{}
  >{\raggedright\arraybackslash}p{(\linewidth - 4\tabcolsep) * \real{0.3333}}
  >{\raggedright\arraybackslash}p{(\linewidth - 4\tabcolsep) * \real{0.3333}}
  >{\raggedright\arraybackslash}p{(\linewidth - 4\tabcolsep) * \real{0.3333}}@{}}
\toprule\noalign{}
\begin{minipage}[b]{\linewidth}\raggedright
时间尺度
\end{minipage} & \begin{minipage}[b]{\linewidth}\raggedright
价格行为
\end{minipage} & \begin{minipage}[b]{\linewidth}\raggedright
分析框架
\end{minipage} \\
\midrule\noalign{}
\endhead
\bottomrule\noalign{}
\endlastfoot
长期 & 价格具有伸缩性，市场趋于出清 & 古典理论（classical theory） \\
短期 & 价格具有黏性，需求可能不等于供给 & 经济周期理论（business cycle theory） \\
\end{longtable}

短期价格黏性可以解释劳动力过剩形成的失业，以及商品需求不足造成的库存积压。长期中，价格充分调整后，经济运行机制可能发生变化。
